\chapter{Daftar Dependensi DecMed}
\label{lampiran:daftar-dependensi-perangkat-lunak}

Daftar ini hanya mencakup dependensi langsung yang dideklarasikan pada masing-masing komponen, bukan seluruh dependensi transitif. Dependensi internal yang ditandai dengan \texttt{path} merupakan \textit{crate} lokal yang dikembangkan sebagai bagian dari sistem. Tabel~\ref{tab:dependensi-umum} sampai Tabel~\ref{tab:dependensi-decmed-macaroon-auth} merinci dependensi tersebut per kelompok komponen.

\begin{longtable}{|>{\raggedright\arraybackslash}p{0.55\textwidth}|>{\raggedright\arraybackslash}p{0.38\textwidth}|}
	\caption{Dependensi Umum yang Digunakan pada Beberapa Komponen}
	\label{tab:dependensi-umum}                                       \\
	\hline
	\textbf{\textit{Crate}}               & \textbf{Versi/Sumber}     \\
	\hline
	\endfirsthead
	\hline
	\textbf{\textit{Crate}}               & \textbf{Versi/Sumber}     \\
	\hline
	\endhead

	\texttt{serde}                        & \texttt{1}                \\ \hline
	\texttt{serde\_json}                  & \texttt{1}                \\ \hline
	\texttt{sha2}                         & \texttt{0.10}             \\ \hline
	\texttt{uuid}                         & \texttt{1}                \\ \hline
	\texttt{bip39}                        & \texttt{2}                \\ \hline
	\texttt{iota-sdk}                     & Git rev. \texttt{3d7c588} \\ \hline
	\texttt{iota-keys}                    & Git rev. \texttt{3d7c588} \\ \hline
	\texttt{iota-config}                  & Git rev. \texttt{3d7c588} \\ \hline
	\texttt{iota-types}                   & Git rev. \texttt{3d7c588} \\ \hline
	\texttt{iota-json-rpc-types}          & Git rev. \texttt{3d7c588} \\ \hline
	\texttt{shared\_crypto}               & Git rev. \texttt{3d7c588} \\ \hline
	\texttt{move-core-types}              & Git rev. \texttt{3d7c588} \\ \hline
	\texttt{dotenv}                       & \texttt{0.15}             \\ \hline
	\texttt{bcs}                          & \texttt{0.1.4}            \\ \hline
	\texttt{schemars}                     & \texttt{0.8}              \\ \hline
	\texttt{rand}                         & \texttt{0.9}              \\ \hline
	\texttt{bincode}                      & \texttt{2}                \\ \hline
	\texttt{umbral-pre}                   & \texttt{0.11}             \\ \hline
	\texttt{aes-gcm}                      & \texttt{0.10}             \\ \hline
	\texttt{argon2}                       & \texttt{0.5.3}            \\ \hline
	\texttt{hex}                          & \texttt{0.4.3}            \\ \hline
	\texttt{regex}                        & \texttt{1.11.1}           \\ \hline
	\texttt{generic-array}                & \texttt{1.2.0}            \\ \hline
	\texttt{k256}                         & \texttt{0.13.4}           \\ \hline
	\texttt{elliptic-curve}               & \texttt{0.13.8}           \\ \hline
	\texttt{chrono}                       & \texttt{0.4}              \\ \hline
	\texttt{base64}                       & \texttt{0.22}             \\ \hline
	\texttt{strum}/\texttt{strum\_macros} & \texttt{0.27}             \\ \hline
	\texttt{anyhow}                       & \texttt{1.0}              \\ \hline
	\texttt{thiserror}                    & \texttt{2.0}              \\ \hline
\end{longtable}

Tabel~\ref{tab:dependensi-umum} menunjukkan fondasi bersama untuk serialisasi, kriptografi, dan integrasi IOTA yang dipakai lintas komponen.

\begin{longtable}{|>{\raggedright\arraybackslash}p{0.55\textwidth}|>{\raggedright\arraybackslash}p{0.38\textwidth}|}
	\caption{Dependensi PRE}
	\label{tab:dependensi-backend-pre}                    \\
	\hline
	\textbf{\textit{Crate}}       & \textbf{Versi/Sumber} \\
	\hline
	\endfirsthead
	\hline
	\textbf{\textit{Crate}}       & \textbf{Versi/Sumber} \\
	\hline
	\endhead

	\texttt{axum}                 & \texttt{0.8.4}        \\ \hline
	\texttt{axum-extra}           & \texttt{0.10.1}       \\ \hline
	\texttt{tower}                & \texttt{0.5.2}        \\ \hline
	\texttt{tokio}                & \texttt{1.44.2}       \\ \hline
	\texttt{reqwest}              & \texttt{0.12}         \\ \hline
	\texttt{redis}                & \texttt{0.32}         \\ \hline
	\texttt{r2d2}                 & \texttt{0.8}          \\ \hline
	\texttt{dotenvy}              & \texttt{0.15.7}       \\ \hline
	\texttt{decmed-rme-segment}   & \texttt{path}         \\ \hline
	\texttt{decmed-macaroon-auth} & \texttt{path}         \\ \hline
\end{longtable}

Tabel~\ref{tab:dependensi-backend-pre} menegaskan bahwa \textit{Server} PRE bertumpu pada \textit{stack} layanan web asinkron, Redis, serta \textit{crate} lokal segmentasi dan otorisasi.

\begin{longtable}{|>{\raggedright\arraybackslash}p{0.55\textwidth}|>{\raggedright\arraybackslash}p{0.38\textwidth}|}
	\caption{Dependensi \textit{Client}}
	\label{tab:dependensi-client-tauri}                       \\
	\hline
	\textbf{\textit{Crate}}           & \textbf{Versi/Sumber} \\
	\hline
	\endfirsthead
	\hline
	\textbf{\textit{Crate}}           & \textbf{Versi/Sumber} \\
	\hline
	\endhead

	\texttt{tauri}                    & \texttt{2}            \\ \hline
	\texttt{tauri-plugin-opener}      & \texttt{2}            \\ \hline
	\texttt{tauri-plugin-http}        & \texttt{2}            \\ \hline
	\texttt{tauri-plugin-fs}          & \texttt{2}            \\ \hline
	\texttt{keyring}                  & \texttt{3}            \\ \hline
	\texttt{rustls-platform-verifier} & \texttt{0.5}          \\ \hline
	\texttt{rqrr}                     & \texttt{0.9}          \\ \hline
	\texttt{image}                    & \texttt{0.25}         \\ \hline
	\texttt{tokio-stream}             & \texttt{0.1.17}       \\ \hline
	\texttt{decmed-rme-segment}       & \texttt{path}         \\ \hline
	\texttt{decmed-macaroon-auth}     & \texttt{path}         \\ \hline
	\texttt{tauri-build}              & \texttt{2}            \\ \hline
\end{longtable}

Tabel~\ref{tab:dependensi-client-tauri} menunjukkan ketergantungan antarmuka desktop Tauri, penyimpanan kunci lokal, serta kemampuan pemindaian QR pada sisi \textit{client}.

\begin{longtable}{|>{\raggedright\arraybackslash}p{0.55\textwidth}|>{\raggedright\arraybackslash}p{0.38\textwidth}|}
	\caption{Dependensi \texttt{decmed-rme-segment}}
	\label{tab:dependensi-decmed-rme-segment}       \\
	\hline
	\textbf{\textit{Crate}} & \textbf{Versi/Sumber} \\
	\hline
	\endfirsthead
	\hline
	\textbf{\textit{Crate}} & \textbf{Versi/Sumber} \\
	\hline
	\endhead

	\texttt{base64}         & \texttt{0.22}         \\ \hline
	\texttt{serde}          & \texttt{1}            \\ \hline
	\texttt{serde\_json}    & \texttt{1}            \\ \hline
	\texttt{sha2}           & \texttt{0.10}         \\ \hline
	\texttt{uuid}           & \texttt{1}            \\ \hline
\end{longtable}

Tabel~\ref{tab:dependensi-decmed-rme-segment} menegaskan bahwa \textit{crate} segmentasi RME bersifat ringan dan berfokus pada serialisasi serta hashing.

\begin{longtable}{|>{\raggedright\arraybackslash}p{0.55\textwidth}|>{\raggedright\arraybackslash}p{0.38\textwidth}|}
	\caption{Dependensi \texttt{decmed-macaroon-auth}}
	\label{tab:dependensi-decmed-macaroon-auth}         \\
	\hline
	\textbf{\textit{Crate}}     & \textbf{Versi/Sumber} \\
	\hline
	\endfirsthead
	\hline
	\textbf{\textit{Crate}}     & \textbf{Versi/Sumber} \\
	\hline
	\endhead

	\texttt{chrono}             & \texttt{0.4}          \\ \hline
	\texttt{serde}              & \texttt{1}            \\ \hline
	\texttt{serde\_json}        & \texttt{1}            \\ \hline
	\texttt{thiserror}          & \texttt{2.0}          \\ \hline
	\texttt{uuid}               & \texttt{1}            \\ \hline
	\texttt{sha2}               & \texttt{0.10}         \\ \hline
	\texttt{hex}                & \texttt{0.4}          \\ \hline
	\texttt{macaroon}           & \texttt{path}         \\ \hline
	\texttt{decmed-rme-segment} & \texttt{path}         \\ \hline
\end{longtable}

Tabel~\ref{tab:dependensi-decmed-macaroon-auth} menunjukkan bahwa otorisasi Macaroon dibangun di atas pustaka token lokal dan tipe kategori dari \texttt{decmed-rme-segment}.
